\chapter{Activation Functions and Normalization}
\label{chap:activation_normalization}

\begin{tldr}
Activation functions introduce nonlinearity; normalization techniques stabilize training. This chapter covers the evolution of activations (sigmoid → ReLU → GELU → SwiGLU), normalization methods (BatchNorm, LayerNorm, RMSNorm), norm placement (pre- vs.\ post-norm and the 2024--26 hybrids), and attention-logit control (QK-norm, soft-capping). Understanding why these components work—and their failure modes—is essential for debugging training issues, from a dead ReLU on one GPU to a loss spike at 100B tokens.
\end{tldr}

Activation functions and normalization are building blocks interviewers use to test debugging intuition. Expect questions about vanishing gradients from sigmoid, why transformers use LayerNorm instead of BatchNorm, dead ReLU neurons, and—at L6 and above—how normalization placement and attention-logit control keep large training runs from diverging.

\section{Activation Functions}
\label{sec:activations}

\subsection{Why Nonlinearity?}

Without activation functions, deep networks collapse to linear models:
\begin{equation}
f(\vx) = \mW_L \cdots \mW_2 \mW_1 \vx = \mW_{eff} \vx
\end{equation}

Nonlinearity enables learning complex decision boundaries and representations.

\subsection{Classical Activations}

\subsubsection{Sigmoid}

\begin{mathresult}{Sigmoid}
\begin{equation}
\sigma(x) = \frac{1}{1 + e^{-x}}, \quad \sigma'(x) = \sigma(x)(1 - \sigma(x))
\end{equation}
\end{mathresult}

\textbf{Properties}:
\begin{itemize}
    \item Output range: $(0, 1)$
    \item Smooth, differentiable everywhere
    \item Maximum gradient: 0.25 (at $x=0$)
\end{itemize}

\textbf{Problems}:
\begin{itemize}
    \item \textbf{Vanishing gradients}: For $|x| > 4$, gradient $\approx 0$
    \item \textbf{Not zero-centered}: Outputs always positive, causes zig-zag gradients
    \item \textbf{Expensive}: Exponential computation
\end{itemize}

\subsubsection{Tanh}

\begin{equation}
\tanh(x) = \frac{e^x - e^{-x}}{e^x + e^{-x}} = 2\sigma(2x) - 1
\end{equation}

\textbf{Improvement over sigmoid}: Zero-centered output $(-1, 1)$

\textbf{Still has}: Vanishing gradients for $|x| > 2$

\subsection{ReLU Family}

\subsubsection{ReLU (Rectified Linear Unit)}

\begin{mathresult}{ReLU}
\begin{equation}
\text{ReLU}(x) = \max(0, x), \quad \text{ReLU}'(x) = \begin{cases} 1 & x > 0 \\ 0 & x \leq 0 \end{cases}
\end{equation}
\end{mathresult}

\textbf{Advantages}:
\begin{itemize}
    \item No vanishing gradient for $x > 0$
    \item Sparse activations (biological plausibility)
    \item Computationally efficient
    \item Enables much deeper networks
\end{itemize}

\textbf{Problems}:
\begin{itemize}
    \item \textbf{Dying ReLU}: Neurons with negative input have zero gradient, never recover
    \item \textbf{Not zero-centered}: All positive outputs
\end{itemize}

\subsubsection{Leaky ReLU and PReLU}

\begin{align}
\text{LeakyReLU}(x) &= \max(\alpha x, x), \quad \alpha = 0.01 \text{ (typically)} \\
\text{PReLU}(x) &= \max(\alpha x, x), \quad \alpha \text{ learned}
\end{align}

\textbf{Fixes dying ReLU}: Small gradient for negative inputs.

\subsubsection{ELU (Exponential Linear Unit)}

\begin{equation}
\text{ELU}(x) = \begin{cases} x & x > 0 \\ \alpha(e^x - 1) & x \leq 0 \end{cases}
\end{equation}

\textbf{Properties}:
\begin{itemize}
    \item Smooth at $x = 0$
    \item Mean activations closer to zero
    \item Saturates for negative values (noise robustness)
\end{itemize}

%--- Q6: Dying ReLU (after ReLU Family section) ---

\begin{interviewq}
\textbf{Q: Explain the dying ReLU problem. How do modern activations solve it?}

\badgeLfive{L5} \quad \badgetype{Conceptual} \quad \badgefreq{\freqhigh}

\stronganswer

\textbf{The dying ReLU problem:} ReLU outputs zero for all negative inputs: $\text{ReLU}(x) = \max(0,x)$. Critically, its gradient is also zero for $x \leq 0$. If a neuron's pre-activation is negative for \textit{every} sample in the dataset, it receives zero gradient and its weights never update---the neuron is permanently ``dead.''

\textbf{How neurons die:} (1) A large gradient update pushes the bias to a very negative value, shifting the pre-activation distribution below zero. (2) Poor initialization (e.g., large negative biases) starts neurons in the dead zone. (3) High learning rates amplify both effects. Once dead, the neuron can never recover because it receives exactly zero gradient---there is no mechanism to ``revive'' it. In deep networks, this can cascade: dead neurons in layer $l$ reduce the expressivity available to layer $l+1$, potentially killing more neurons.

\textbf{Diagnosis:} Compute the fraction of zero activations per layer. If $>$30--40\% of neurons are consistently zero across the dataset, dying ReLU is a problem. In PyTorch: register a forward hook that logs \texttt{(output == 0).float().mean()}.

\textbf{Modern solutions, from simplest to most complex:}

\textit{Leaky ReLU / PReLU:} $f(x) = \max(\alpha x, x)$ with $\alpha = 0.01$ (Leaky) or learned (PReLU). The small negative slope ensures a nonzero gradient for $x < 0$, so neurons can always receive updates. Simple, effective, but the negative slope is somewhat arbitrary.

\textit{ELU:} $f(x) = x$ for $x > 0$, $\alpha(e^x - 1)$ for $x \leq 0$. Smooth at zero, and the negative saturation pushes mean activations closer to zero. But the exponential is more expensive than ReLU.

\textit{GELU:} $f(x) = x \cdot \Phi(x)$. Smoothly gates the input by its Gaussian percentile. For large negative inputs, the output and gradient are small but nonzero. For values near zero, the transition is smooth, avoiding the gradient discontinuity. This ``soft thresholding'' is the key conceptual advance: instead of a hard zero/pass decision, GELU makes a probabilistic decision.

\textit{SiLU/Swish:} $f(x) = x \cdot \sigma(x)$. Self-gated: the sigmoid of $x$ gates $x$ itself. Has a small negative region (minimum $\approx -0.28$ at $x \approx -1.28$), which empirically acts as regularization. Used in EfficientNet and as part of SwiGLU.

All modern activations solve dying ReLU through the same principle: ensuring \textbf{nonzero gradients for negative inputs}. The evolution is from simple fixes (Leaky ReLU) to smooth, expressive functions (GELU, SiLU) that additionally provide better optimization landscapes.

\redflags
\begin{itemize}
    \item ``Dying ReLU happens when the output is zero''---the problem is that the \textit{gradient} is zero, preventing recovery.
    \item ``Just use a smaller learning rate''---this reduces the probability but doesn't eliminate the fundamental issue.
    \item Not knowing that GELU still allows near-zero activations for negative inputs; the key is that the gradient is nonzero.
\end{itemize}

\interviewtest{Understanding of the interaction between activation functions and gradient flow, the ability to explain \textit{why} a neuron can never recover (zero gradient = zero update), and knowledge of how each modern activation addresses this.}

\goodgreat
\begin{itemize}
    \item \badgeLfive{L5} Explains the problem clearly (zero gradient for negative inputs prevents recovery), names 2--3 solutions, and can diagnose it in practice.
    \item \badgeLsix{L6} Connects to initialization theory (He initialization was designed to keep pre-activations balanced around zero for ReLU), explains the cascade effect in deep networks, and discusses the interaction with normalization (LayerNorm re-centers activations, reducing the probability of dying neurons).
    \item \badgeLseven{L7} Discusses the representation capacity implications---dead neurons reduce the effective width of the network. Connects to the lottery ticket hypothesis (sparse networks with many ``dead'' neurons can still work if the right neurons survive).
\end{itemize}

\followups{``How does He initialization specifically address dying ReLU?'' ``Can a dead ReLU neuron ever recover?'' ``What is the relationship between activation sparsity and model quality?''}
\end{interviewq}

\subsection{Modern Activations}

\subsubsection{GELU (Gaussian Error Linear Unit)}

\begin{mathresult}{GELU}
\begin{equation}
\text{GELU}(x) = x \cdot \Phi(x) = x \cdot P(Z \leq x), \quad Z \sim \mathcal{N}(0, 1)
\end{equation}

Approximation:
\begin{equation}
\text{GELU}(x) \approx 0.5x\left(1 + \tanh\left(\sqrt{2/\pi}(x + 0.044715x^3)\right)\right)
\end{equation}
\end{mathresult}

\textbf{Intuition}: Weight input by its percentile under Gaussian distribution.

\textbf{Used in}: BERT, GPT-2, GPT-3, ViT

\subsubsection{SiLU / Swish}

\begin{equation}
\text{SiLU}(x) = x \cdot \sigma(x) = \frac{x}{1 + e^{-x}}
\end{equation}

\textbf{Properties}:
\begin{itemize}
    \item Self-gated: Uses input to gate itself
    \item Smooth, non-monotonic
    \item Slight negative region (regularization effect)
\end{itemize}

\textbf{Used in}: EfficientNet, LLaMA (in SwiGLU)

\subsubsection{SwiGLU (Swish-Gated Linear Unit)}

\begin{mathresult}{SwiGLU}
For FFN in Transformers:
\begin{equation}
\text{SwiGLU}(\vx, \mW, \mV, \mW_2) = (\text{Swish}(\vx\mW) \odot \vx\mV) \mW_2
\end{equation}
\end{mathresult}

\textbf{Key insight}: Multiplicative gating provides richer interactions.

\textbf{Used in}: LLaMA, PaLM, modern LLMs

\textbf{Note}: Requires 3 weight matrices instead of 2, so FFN expansion factor reduced (8/3 instead of 4).

\textbf{Concrete numbers (Llama-2-7B, $d = 4096$)}: the SwiGLU hidden dimension is $\frac{8}{3} \times 4096 \approx 10923$, rounded to $11008$ for hardware alignment. FFN parameters per layer: $3 \times 4096 \times 11008 \approx 135$M, so the 32 FFN blocks hold $\approx 4.3$B of the model's 6.7B parameters---about two thirds. Activation memory is similarly FFN-dominated: a transformer layer stores roughly $34\,d$ bytes per token for the backward pass in mixed precision (Korthikanti et al.), i.e., $\approx 140$~KB per token per layer at $d = 4096$ excluding the attention score matrices (which FlashAttention avoids storing)---and the widest tensors in that budget are the three $11008$-dim SwiGLU intermediates. This is why activation checkpointing targets the FFN first.

%--- Q3: GELU vs ReLU vs SwiGLU (after Modern Activations) ---

\begin{interviewq}
\textbf{Q: Why GELU over ReLU? Why SwiGLU over GELU?}

\badgeLfive{L5} \quad \badgetype{Trade-off} \quad \badgefreq{\freqhigh}

\stronganswer

\textbf{ReLU to GELU.} ReLU applies a hard threshold: $\text{ReLU}(x) = \max(0,x)$. This has two problems for Transformers. (1) The \textit{dying ReLU} problem: neurons receiving consistently negative inputs have exactly zero gradient and never recover. In deep Transformers with hundreds of sublayers, this compounds---entire channels can go permanently dead. (2) The \textit{hard discontinuity} at $x=0$ means the gradient is undefined at zero and jumps from 0 to 1, creating a non-smooth loss landscape.

GELU smooths this by weighting the input by its Gaussian CDF: $\text{GELU}(x) = x \cdot \Phi(x)$. Near zero, GELU has a smooth transition instead of a hard cutoff. This means small negative inputs get small but \textit{nonzero} outputs and gradients, preventing dead neurons. The smooth curvature also benefits optimization---Adam and other adaptive optimizers work better with smooth gradients. Empirically, GELU gives 0.5--1\% improvement on Transformer benchmarks (BERT, GPT-2, ViT) with negligible compute overhead.

\textbf{GELU to SwiGLU.} SwiGLU replaces the two-matrix FFN ($\text{FFN}(x) = \text{GELU}(x\mW_1)\mW_2$) with a three-matrix gated structure: $\text{SwiGLU}(x) = (\text{Swish}(x\mW_1) \odot x\mV)\mW_2$. The key innovation is the \textit{multiplicative gating}: $x\mV$ acts as a gate on $\text{Swish}(x\mW_1)$. This element-wise multiplication creates richer feature interactions within the FFN, allowing the network to learn more expressive transformations. The Swish activation ($x\sigma(x)$) is itself smooth and self-gated.

The cost is an extra weight matrix ($\mV$), so to keep parameter count constant, the hidden dimension is reduced from $4d$ to $\frac{8}{3}d$. Despite fewer hidden units, the multiplicative interaction compensates, and SwiGLU consistently outperforms GELU by 0.5--2\% across model scales (shown in PaLM and LLaMA ablations). This is why essentially all modern LLMs (LLaMA, Mistral, Gemma, Qwen) use SwiGLU.

\redflags
\begin{itemize}
    \item ``GELU is just a smoother version of ReLU''---this undersells the gradient flow and dying neuron benefits.
    \item ``SwiGLU is better because it's newer''---need to explain the gating mechanism.
    \item Not mentioning the parameter budget trade-off (3 matrices instead of 2, reduced hidden dim).
\end{itemize}

\interviewtest{Whether you understand the progression of design choices, the specific failure modes each addresses, and the engineering trade-offs (compute, parameters, expressivity) at each step.}

\goodgreat
\begin{itemize}
    \item \badgeLfive{L5} Clearly explains dying ReLU fix (GELU) and gating mechanism (SwiGLU), mentions the $\frac{8}{3}d$ hidden dim.
    \item \badgeLsix{L6} Cites ablation results from PaLM/LLaMA papers, explains why multiplicative interactions increase expressivity (bilinear maps can represent higher-order feature interactions), discusses compute-quality Pareto frontier.
    \item \badgeLseven{L7} Connects to the broader GLU family (original GLU by Dauphin et al., ReGLU, GEGLU), discusses how activation choice interacts with scaling laws, and notes that the benefit of SwiGLU is more pronounced at larger model scales.
\end{itemize}

\followups{``What is the GLU family and how do variants compare?'' ``How does activation choice interact with initialization?'' ``Would you use SwiGLU in a CNN?''}
\end{interviewq}

\subsection{Activation Selection Guide}

\begin{table}[H]
\centering
\caption{Activation function selection}
\begin{tabularx}{\textwidth}{lXX}
\toprule
\textbf{Architecture} & \textbf{Recommended} & \textbf{Alternative} \\
\midrule
CNN (general) & ReLU & LeakyReLU \\
BERT-style encoder & GELU & ReLU \\
GPT-style decoder & SwiGLU (or GEGLU) & GELU (legacy, GPT-2 era) \\
Modern LLMs & SwiGLU & GELU \\
Output (classification) & Softmax & Sigmoid (multi-label) \\
Output (regression) & None (linear) & ReLU (non-negative) \\
\bottomrule
\end{tabularx}
\end{table}

%==============================================================================
\section{Normalization Techniques}
\label{sec:normalization}
%==============================================================================

\subsection{Why Normalization?}

\textbf{The historical story (contested): internal covariate shift.} BatchNorm was originally motivated by \textit{internal covariate shift} (Ioffe and Szegedy, 2015)---the idea that layer-input distributions shift as earlier layers update, making optimization harder. Treat this as history, not mechanism: Santurkar et al.\ (2018) showed BatchNorm helps even when internal covariate shift is artificially \textit{increased}, so ICS cannot be the explanation.

\textbf{The current understanding: loss-landscape smoothing.} Normalization improves the Lipschitzness of the loss and its gradients---gradients become more predictive, so larger learning rates are stable and training is less sensitive to initialization. In interviews, reciting the ICS story uncritically is itself a red flag; ``is the internal-covariate-shift explanation actually true?'' is a deliberate follow-up.

\textbf{Normalization benefits}:
\begin{itemize}
    \item Stabilizes training (smoother loss landscape, more predictive gradients)
    \item Enables higher learning rates
    \item Reduces sensitivity to initialization
    \item Acts as regularization
\end{itemize}

\subsection{Batch Normalization}

\begin{mathresult}{Batch Normalization}
For a mini-batch $\mathcal{B} = \{x_1, \ldots, x_m\}$:
\begin{align}
\mu_\mathcal{B} &= \frac{1}{m}\sum_{i=1}^m x_i \\
\sigma_\mathcal{B}^2 &= \frac{1}{m}\sum_{i=1}^m (x_i - \mu_\mathcal{B})^2 \\
\hat{x}_i &= \frac{x_i - \mu_\mathcal{B}}{\sqrt{\sigma_\mathcal{B}^2 + \epsilon}} \\
y_i &= \gamma \hat{x}_i + \beta
\end{align}
where $\gamma, \beta$ are learned parameters.
\end{mathresult}

\textbf{Training vs. Inference}:
\begin{itemize}
    \item Training: Use batch statistics
    \item Inference: Use running mean/variance from training
\end{itemize}

\textbf{Limitations}:
\begin{itemize}
    \item Requires sufficiently large batch sizes
    \item Different behavior train vs. test
    \item Problematic for sequence models (variable length)
    \item Not used in Transformers
\end{itemize}

%--- Q5: Train vs inference normalization (after BatchNorm) ---

\begin{interviewq}
\textbf{Q: Your model shows different behavior at training vs inference. What normalization issue could cause this?}

\badgeLsix{L6} \quad \badgetype{Debugging} \quad \badgefreq{\freqmed}

\stronganswer

The most common culprit is \textbf{BatchNorm's train/eval mode discrepancy}. During training, BatchNorm uses per-batch statistics ($\mu_\mathcal{B}, \sigma^2_\mathcal{B}$). During inference, it uses exponential moving averages ($\mu_{running}, \sigma^2_{running}$) accumulated during training. Several things can go wrong.

\textbf{Scenario 1: Forgot to call model.eval().} This is the most common cause. In PyTorch, if you don't switch to eval mode, BatchNorm continues using batch statistics at inference time. With batch size = 1 (common during inference), these statistics are just the single sample's mean and variance, leading to incorrect normalization. The fix is trivial but the bug is insidious because the model may appear to work but with degraded quality. Verify with the \texttt{model.training} attribute, and check the mirror-image bug too: if the running mean and variance are all zeros/ones, the model was in eval mode \textit{during training} and the running statistics were never accumulated. A quick behavioral diagnostic: run inference with batch sizes 1, 16, and 64---if the outputs change with batch size, BatchNorm is computing live statistics (LayerNorm would be batch-size-invariant).

\textbf{Scenario 2: Running statistics diverged from true test distribution.} If training data distribution differs from test data (domain shift), running statistics are wrong. This is common when: (a) training uses heavy data augmentation that shifts activation distributions, (b) the model is fine-tuned on a different domain without resetting running statistics, (c) training used dropout, which changes activation magnitudes (BatchNorm after dropout computes statistics on dropped-out activations during training but full activations at inference). Two things to check and fix: the momentum parameter (PyTorch's default \texttt{momentum=0.1} means running stats are an exponential average with 0.1 weight on each new batch---too high and they track the last batches, too low and they lag), and, for the augmentation case, recompute running statistics in a final pass over the training set with augmentation disabled; comparing these fresh statistics against the accumulated ones is itself a diagnostic.

\textbf{Scenario 3: Multi-GPU training without SyncBatchNorm.} Each GPU computes BatchNorm statistics from its local mini-batch. If the per-GPU batch size is small (e.g., 4), these statistics are noisy. The running averages accumulate noisy estimates, which diverge from the true population statistics. At inference, these inaccurate running statistics are used, degrading performance. The fix is to use \texttt{torch.nn.SyncBatchNorm} or switch to GroupNorm/LayerNorm.

\textbf{Scenario 4: Batch composition at inference differs from training.} If the inference batch mixes different data types (e.g., images of different resolutions padded to the same size, or sequences of wildly different lengths), and BatchNorm is applied before masking, the statistics are polluted by padding tokens or resized artifacts.

\textbf{Why LayerNorm/RMSNorm avoid all of this:} They normalize within each sample independently, so train/eval behavior is identical. No running statistics, no batch-size sensitivity, no multi-GPU synchronization needed. This is a key (often underappreciated) reason Transformers use LayerNorm.

\redflags
\begin{itemize}
    \item ``Just use LayerNorm instead''---this is the right eventual recommendation but the interviewer wants you to diagnose the BatchNorm issue first.
    \item Not mentioning model.eval()---this is the \#1 real-world cause.
    \item Confusing BatchNorm's train/eval behavior with dropout's train/eval behavior.
\end{itemize}

\interviewtest{Deep understanding of how BatchNorm behaves differently at train vs. inference time, the specific mechanisms (running statistics vs. batch statistics), and practical debugging skills for a class of bugs that are notoriously hard to diagnose because the model ``mostly works.''}

\goodgreat
\begin{itemize}
    \item \badgeLfive{L5} Identifies model.eval() and running statistics issues, knows LayerNorm avoids the problem.
    \item \badgeLsix{L6} Adds multi-GPU SyncBatchNorm, dropout-before-BatchNorm interaction, and domain shift scenarios. Proposes concrete diagnostics: the batch-size-sensitivity test (outputs at batch 1 vs.\ 16 vs.\ 64), inspecting running statistics for never-updated (all-zero) values, and recomputing fresh statistics on un-augmented data for comparison.
    \item \badgeLseven{L7} Discusses the deeper issue that any normalization depending on batch composition creates an implicit dependency between samples, which violates the assumption of independent inference. Proposes normalization-robust evaluation pipelines in production: always log running statistics, and add automated batch-size-invariance tests in CI so this entire bug class is caught mechanically. Connects to the broader design principle of using sample-independent normalization in production systems.
\end{itemize}

\followups{``How would you verify that running statistics are correct?'' ``Does dropout cause a similar train/test discrepancy?'' ``How does BatchNorm interact with gradient accumulation?'' ``How would you automate detection of this class of bugs?''}
\end{interviewq}

\subsection{Layer Normalization}

\begin{mathresult}{Layer Normalization}
For each sample, normalize across features:
\begin{align}
\mu &= \frac{1}{d}\sum_{i=1}^d x_i \\
\sigma^2 &= \frac{1}{d}\sum_{i=1}^d (x_i - \mu)^2 \\
\hat{x}_i &= \frac{x_i - \mu}{\sqrt{\sigma^2 + \epsilon}} \\
y_i &= \gamma_i \hat{x}_i + \beta_i
\end{align}
\end{mathresult}

\textbf{Key difference from BatchNorm}: Normalizes across features (within sample), not across batch.

\textbf{Advantages}:
\begin{itemize}
    \item Works with any batch size (even 1)
    \item Same behavior train and test
    \item Natural for sequence models
\end{itemize}

\textbf{Used in}: Transformers (BERT, GPT, etc.)

%--- Q1: BatchNorm-Transformer incompatibility (after LayerNorm; merged with former Q7) ---

\begin{interviewq}
\textbf{Q: Why don't Transformers use BatchNorm? Derive the reasoning from first principles.}

\badgeLsix{L6} \quad \badgetype{First Principles} \quad \badgefreq{\freqmed}

\stronganswer

Derive this from the mechanics of both BatchNorm and the Transformer architecture, rather than asserting the conclusion.

\textbf{BatchNorm formulation.} For input tensor $\vx \in \R^{B \times T \times d}$ (batch, sequence length, features), BatchNorm computes per-feature statistics across the batch and sequence dimensions:
\begin{equation*}
\mu_j = \frac{1}{BT}\sum_{b=1}^B\sum_{t=1}^T x_{b,t,j}, \quad \sigma_j^2 = \frac{1}{BT}\sum_{b=1}^B\sum_{t=1}^T (x_{b,t,j} - \mu_j)^2
\end{equation*}
Then normalizes: $\hat{x}_{b,t,j} = (x_{b,t,j} - \mu_j)/\sqrt{\sigma_j^2 + \epsilon}$.

\textbf{Problem 1: Token heterogeneity---the batch dimension means something different than in CNNs.} In a CNN, feature $j$ at spatial position $(h,w)$ across different images measures a consistent quantity (e.g., edge presence), so averaging across the batch is meaningful. In a Transformer, tokens in different positions play different roles (subject, verb, beginning-of-sequence marker, padding): feature $j$ at $t=1$ and feature $j$ at $t=100$ measure different things. Averaging across both batch and time conflates semantically unrelated activations---the statistics $\mu_j$ and $\sigma_j^2$ are averages over a mixture of distributions, not a single coherent distribution.

\textbf{Problem 2: Padding corruption.} With variable-length sequences, shorter sequences are padded. If padding tokens are included in the BatchNorm computation, they corrupt the statistics. If excluded (masked BatchNorm), the effective sample count depends on how much padding each batch contains, so the statistics stay noisy; a per-position variant (separate statistics per $t$) is worse still, since positions near the maximum length see very few real tokens.

\textbf{Problem 3: Autoregressive inference.} At generation time, the model processes one token at a time with batch size 1. BatchNorm must use running statistics from training. But training statistics were computed over \textit{full sequences} (including future tokens), while inference sees only \textit{partial sequences}. The activation distribution conditioned on ``sequence so far'' differs from the training distribution conditioned on ``all tokens.'' This is a fundamental distribution shift, on top of the ordinary train/eval discrepancy BatchNorm already carries.

\textbf{Problem 4: Gradient coupling.} BatchNorm creates a dependency between samples in the batch through the shared statistics. The gradient $\partial \loss / \partial x_{b,t,j}$ depends on all other samples in the batch via $\mu_j$ and $\sigma_j^2$. Be precise here: this coupling is \textit{not} Transformer-specific---CNNs with BatchNorm have exactly the same property. It is a general cost of batch-statistic normalization (noisier gradients, batch-composition dependence, samples no longer independent at inference); the honest framing is that Transformers have no reason to pay this cost, since per-token normalization works at least as well for them.

\textbf{LayerNorm resolves all four issues} by computing statistics across features within each token: $\mu_{b,t} = \frac{1}{d}\sum_{j=1}^d x_{b,t,j}$. Each token is normalized independently: no batch dependence, no padding issues, no train/inference mismatch, no gradient coupling between samples. This also aligns with how attention processes sequences---each token's representation is updated independently through Q/K/V projections before interacting with other tokens.

\redflags
\begin{itemize}
    \item ``BatchNorm just doesn't work for sequences''---hand-waving; the question asks you to derive the mechanism.
    \item Missing the autoregressive inference argument (the strongest theoretical objection).
    \item Not writing out the BatchNorm statistics formula to show where the batch/time dimensions interact, or not distinguishing the batch-dimension semantics between CNNs and Transformers.
\end{itemize}

\interviewtest{Whether you can reason from the mathematical formulation of a normalization method to its architectural consequences---deriving \textit{why} the choice was made rather than pattern-matching ``Transformers use LayerNorm.''}

\goodgreat
\begin{itemize}
    \item \badgeLfive{L5} States 2--3 problems clearly, including variable lengths and the train/test discrepancy.
    \item \badgeLsix{L6} Writes out the formulas showing how batch and time dimensions interact in BatchNorm statistics, contrasts with the LayerNorm formula, explains why the batch-dimension semantics differ between CNNs and Transformers, addresses gradient coupling honestly, and mentions that PowerNorm and other BatchNorm variants have been tried in Transformers with limited success.
    \item \badgeLseven{L7} Connects to the broader design principle that normalization should be \textit{per-example} where batch composition is variable and should align with the architecture's inductive biases; discusses the loss-landscape-smoothing view of why normalization helps at all; and mentions that some recent work (e.g., nGPT) removes normalization entirely in favor of representation learning on the hypersphere.
\end{itemize}

\followups{``Why RMSNorm over LayerNorm in modern LLMs?'' ``Where is BatchNorm still preferred and why?'' ``Could you make BatchNorm work if you computed statistics per-position?'' ``What is pre-norm vs. post-norm and why does it matter for training stability?''}
\end{interviewq}

\subsection{RMSNorm}

\begin{mathresult}{RMSNorm}
\begin{equation}
\text{RMSNorm}(\vx) = \frac{\vx}{\text{RMS}(\vx)} \odot \bm{\gamma}, \quad \text{RMS}(\vx) = \sqrt{\frac{1}{d}\sum_{i=1}^d x_i^2 + \epsilon}
\end{equation}
\end{mathresult}

\textbf{Simplification of LayerNorm}:
\begin{itemize}
    \item No mean centering (only scale normalization)
    \item No bias term $\beta$
    \item Empirically similar performance
    \item Speed: the norm \textit{op} itself is 10--15\% cheaper, but norm layers are only $\sim$1\% of end-to-end compute on modern fused stacks---the real appeal is simplicity and fewer parameters
\end{itemize}

\textbf{Used in}: LLaMA, Mistral, most modern LLMs

%--- Q4: RMSNorm vs LayerNorm (after RMSNorm) ---

\begin{interviewq}
\textbf{Q: RMSNorm vs LayerNorm---when would you choose each?}

\badgeLfive{L5} \quad \badgetype{Trade-off} \quad \badgefreq{\freqmed}

\stronganswer

LayerNorm normalizes by centering (subtracting mean) and scaling (dividing by standard deviation): $\hat{x}_i = \frac{x_i - \mu}{\sqrt{\sigma^2 + \epsilon}}$, then applies learned affine: $y_i = \gamma_i \hat{x}_i + \beta_i$. RMSNorm drops the mean-centering step entirely: $\hat{x}_i = \frac{x_i}{\text{RMS}(\vx)} \cdot \gamma_i$, where $\text{RMS}(\vx) = \sqrt{\frac{1}{d}\sum_j x_j^2 + \epsilon}$.

\textbf{Choose RMSNorm when:} (1) Training modern LLMs---the normalization op itself is 10--15\% cheaper (no mean computation, no bias parameter), though norm layers account for only $\sim$1\% of end-to-end compute on fused stacks, so the honest motivation is simplicity and fewer parameters rather than material wall-clock savings. (2) The architecture uses pre-norm placement, where the re-centering property of LayerNorm is less critical because the residual stream preserves the mean. (3) You want a simpler implementation (fewer parameters: no $\beta$).

\textbf{Choose LayerNorm when:} (1) Working with encoder-only models like BERT where the original architecture uses LayerNorm and swapping may require re-tuning. (2) Post-norm architectures, where the mean-centering provides additional stability. (3) Tasks where the mean of activations carries important information---removing it may slightly hurt in some settings. (4) When using well-established codebases that default to LayerNorm and switching offers marginal gains.

The empirical insight from Zhang and Sennrich (2019) is that mean-centering contributes negligibly to LayerNorm's effectiveness---the scale normalization (preventing activations from growing or shrinking) does the heavy lifting. This is why RMSNorm achieves comparable quality: it preserves the essential re-scaling while being cheaper to compute.

In practice, the quality difference between LayerNorm and RMSNorm is within noise for most tasks. The choice is driven by simplicity and convention more than by end-to-end speed, and for any new LLM project, RMSNorm is the default recommendation.

\redflags
\begin{itemize}
    \item ``RMSNorm is strictly better''---there are cases where mean-centering matters.
    \item Not knowing that RMSNorm removes both the mean subtraction \textit{and} the bias term $\beta$.
    \item Overstating the speed difference---10--15\% is on the normalization op, not the full model.
\end{itemize}

\interviewtest{Understanding of what each component (mean-centering, variance scaling, learned affine) contributes, and making principled trade-off decisions rather than blindly following trends.}

\goodgreat
\begin{itemize}
    \item \badgeLfive{L5} Knows the formula difference, states RMSNorm is faster with similar quality, and recommends it for new LLM projects.
    \item \badgeLsix{L6} Explains \textit{why} mean-centering is less important (the re-scaling prevents activation explosion, which is the primary benefit), cites the Zhang and Sennrich analysis, and discusses interaction with pre-norm vs. post-norm.
    \item \badgeLseven{L7} Discusses how the choice interacts with residual stream dynamics (pre-norm + RMSNorm means the residual stream is only scaled, not shifted, preserving information in the mean), and mentions emerging approaches like QK-norm and nGPT that further simplify normalization.
\end{itemize}

\followups{``What if you removed normalization entirely?'' ``How does normalization interact with learning rate selection?'' ``What is QK-norm and why is it needed?''}
\end{interviewq}

\subsection{Other Normalization Variants}

\subsubsection{Instance Normalization}

Normalize each channel of each sample independently.
\begin{equation}
\text{IN}(x_{nchw}) = \frac{x_{nchw} - \mu_{nc}}{\sigma_{nc}}
\end{equation}

\textbf{Used in}: Style transfer, image generation

\subsubsection{Group Normalization}

Divide channels into groups, normalize within groups.

\textbf{Used in}: When batch size is very small (object detection)

\subsubsection{Weight Normalization}

Reparameterize weights: $\vw = g \frac{\vv}{\|\vv\|}$

\textbf{Used in}: Some generative models

\subsection{Pre-Norm vs. Post-Norm}
\label{subsec:prenorm_postnorm}

Where the normalization sits relative to the residual connection is one of the most consequential stability decisions in a Transformer:
\begin{align}
\text{Post-norm (original Transformer, BERT):}\quad \vx_{l+1} &= \mathrm{Norm}\big(\vx_l + F_l(\vx_l)\big) \\
\text{Pre-norm (GPT-2 onward, modern default):}\quad \vx_{l+1} &= \vx_l + F_l\big(\mathrm{Norm}(\vx_l)\big)
\end{align}
where $F_l$ is the attention or FFN sublayer.

\begin{figure}[H]
\centering
\begin{tikzpicture}[
    node distance=0.5cm,
    box/.style={rectangle, draw, minimum width=2cm, minimum height=0.6cm, rounded corners, font=\small},
    norm/.style={box, fill=lightyellow},
    sublayer/.style={box, fill=lightblue}
]
    % Post-norm
    \begin{scope}[shift={(0,0)}]
        \node[sublayer] (sub1) {SubLayer};
        \node[above=0.3cm of sub1] (add1) {$+$};
        \node[norm, above=0.3cm of add1] (norm1) {Norm};
        \node[above=0.5cm of norm1] (out1) {};
        \node[below=0.5cm of sub1] (in1) {};

        \draw[->] (in1) -- (sub1);
        \draw[->] (sub1) -- (add1);
        \draw[->] (add1) -- (norm1);
        \draw[->] (norm1) -- (out1);
        \draw[->] (in1.north) -- ++(-0.8,0) |- (add1.west);

        \node[below=1cm of sub1, font=\bfseries] {Post-Norm};
        \node[below=1.4cm of sub1, font=\small] {(Original Transformer)};
    \end{scope}

    % Pre-norm
    \begin{scope}[shift={(5,0)}]
        \node[norm] (norm2) {Norm};
        \node[sublayer, above=0.3cm of norm2] (sub2) {SubLayer};
        \node[above=0.3cm of sub2] (add2) {$+$};
        \node[above=0.5cm of add2] (out2) {};
        \node[below=0.5cm of norm2] (in2) {};

        \draw[->] (in2) -- (norm2);
        \draw[->] (norm2) -- (sub2);
        \draw[->] (sub2) -- (add2);
        \draw[->] (add2) -- (out2);
        \draw[->] (in2.north) -- ++(-0.8,0) |- (add2.west);

        \node[below=1cm of norm2, font=\bfseries] {Pre-Norm};
        \node[below=1.4cm of norm2, font=\small] {(Modern Standard)};
    \end{scope}
\end{tikzpicture}
\caption{Pre-norm vs. post-norm placement in Transformer blocks.}
\end{figure}

\textbf{The gradient argument (Xiong et al., 2020).} In the pre-norm stack the residual path is never normalized, so the Jacobian between any two depths contains a clean identity term:
\begin{equation}
\frac{\partial \vx_L}{\partial \vx_l} = \prod_{k=l}^{L-1}\left(\mI + \frac{\partial F_k(\mathrm{Norm}(\vx_k))}{\partial \vx_k}\right) = \mI + (\text{cross terms}),
\end{equation}
and the loss gradient reaches every layer with an $O(1)$ component regardless of depth. In the post-norm stack, every factor in the product is wrapped in a $\mathrm{Norm}$ Jacobian---no path from the loss to layer $l$ avoids the $L - l$ intervening normalizations, each of which rescales by roughly the inverse of the activation norm. Xiong et al.\ analyzed both at initialization: post-norm concentrates large parameter gradients in the layers near the output while starving the lower layers, and the imbalance grows with depth. This is precisely why the original Transformer recipe depends on learning-rate warmup---at full learning rate, the first few updates to the top post-norm layers are large enough to destroy the model before training settles---and why deep post-norm stacks (18+ layers) historically failed to train without tricks like DeepNet's rescaled residuals. Pre-norm removes most of this warmup dependence, but not all of it: every modern LLM recipe still warms up, just far less delicately.

\textbf{Pre-norm's own pathology: residual-stream growth.} The identity path means each block \emph{adds} to a stream that is never re-normalized, so $\|\vx_l\|$ grows with depth---roughly $\sqrt{l}$ at initialization (a sum of $l$ roughly independent block outputs), and often faster during training. Each block's $O(1)$-normalized contribution is therefore added to an ever-larger stream: later blocks change the stream's direction less and less, and a deep pre-norm network behaves \emph{effectively shallower} than its layer count. This is one reason doubling depth does not double quality in pre-norm models, and why the final norm before the output head is essential (without it, the head sees activations whose scale depends on depth).

\textbf{2024--26 hybrids.} Recent models re-introduce controlled post-norm elements without breaking the identity path. Gemma~2 uses \emph{sandwich norm}: RMSNorm on the sublayer input \emph{and} on its output before the residual addition, $\vx + \mathrm{Norm}(F(\mathrm{Norm}(\vx)))$, capping how much any single block injects into the stream. OLMo~2 instead \emph{reorders} the norm to the sublayer output only, $\vx + \mathrm{Norm}(F(\vx))$---a post-norm variant that leaves the residual path untouched and measurably improved its training stability. \emph{Peri-LN} (2025) formalizes the pattern---normalize both the input and the output of every module---showing it keeps residual-stream variance growth linear rather than exponential while keeping gradient scales even across depth.

%--- Q9: Pre-norm vs post-norm gradient flow (after Pre-Norm vs. Post-Norm) ---

\begin{interviewq}
\textbf{Q: Pre-norm vs. post-norm---derive the gradient-flow difference and explain what modern hybrids fix.}

\badgeLsix{L6} \quad \badgetype{Trade-off} \quad \badgefreq{\freqmed}

\stronganswer

\textbf{Set up the two update rules.} Post-norm: $\vx_{l+1} = \mathrm{Norm}(\vx_l + F_l(\vx_l))$. Pre-norm: $\vx_{l+1} = \vx_l + F_l(\mathrm{Norm}(\vx_l))$. The entire stability story follows from differentiating these.

\textbf{Derive the backward path.} For pre-norm, the chain rule gives
\begin{equation*}
\frac{\partial \loss}{\partial \vx_l} = \frac{\partial \loss}{\partial \vx_L}\prod_{k=l}^{L-1}\big(\mI + \mJ_k\big), \qquad \mJ_k = \frac{\partial F_k(\mathrm{Norm}(\vx_k))}{\partial \vx_k}.
\end{equation*}
Expanding the product, one term is exactly $\mI$: the loss gradient reaches layer $l$ unattenuated, no matter how deep the stack. Everything else is a correction on top of that identity term. For post-norm, each factor is $\mathrm{Norm}'(\cdot)\,(\mI + \partial F_k/\partial \vx_k)$---the norm Jacobian multiplies \textit{every} path, including the residual one, so no identity term survives. A LayerNorm Jacobian scales like the inverse of its input norm, so the product of $L-l$ of them neither preserves gradient magnitude nor treats layers evenly. Xiong et al.\ (2020) make this precise at initialization: post-norm puts large parameter gradients in the layers near the output, the layer-to-layer imbalance grows with depth, and a full-size learning rate destroys the top layers in the first few steps---which is exactly why post-norm training needs warmup (BERT's recipe) and why removing warmup from a post-norm Transformer diverges. Pre-norm's gradients are well-behaved at any depth, so warmup sensitivity drops sharply (though modern recipes still warm up).

\textbf{What pre-norm gives up.} The unnormalized residual stream grows with depth ($\approx\sqrt{l}$ at initialization, often faster in training), so later blocks contribute relatively less and the network behaves effectively shallower than its layer count. The growing stream also feeds ever-larger activations into attention and the output head, which interacts badly with low precision. Post-norm, when it trains at all, avoids this by re-normalizing the stream every block---one reason well-tuned post-norm models can match or beat pre-norm at convergence.

\textbf{What the hybrids fix.} They buy back post-norm's activation control without giving up the identity path. Gemma~2's sandwich norm normalizes each sublayer's input and its output before the residual addition, bounding what any block injects into the stream. OLMo~2 normalizes the sublayer output only ($\vx + \mathrm{Norm}(F(\vx))$), which its ablations showed reduced gradient-norm spikes. Peri-LN unifies these: normalizing into and out of every module keeps residual variance growth linear and gradient scales even across layers. In all three, the second norm sits on the \textit{sublayer output}, never on the residual sum---normalizing the sum is what would reintroduce post-norm's broken identity path.

\redflags
\begin{itemize}
    \item ``Pre-norm is better, period''---without the gradient mechanism, this is memorized trivia, and it cannot explain why Gemma~2 and OLMo~2 deliberately added post-norm elements back.
    \item Claiming pre-norm eliminates the need for warmup entirely---it reduces warmup sensitivity; every modern large-model recipe still warms up.
    \item Placing sandwich norm's second norm on the residual sum rather than the sublayer output---that would break the identity path and defeat the design.
\end{itemize}

\interviewtest{Whether you can convert an architecture diagram into a Jacobian statement and reason about what survives a product over depth---and whether your knowledge extends past ``pre-norm is standard'' to why the 2024--26 generation partially reversed it.}

\goodgreat
\begin{itemize}
    \item \badgeLfive{L5} Knows the two placements, that pre-norm is the modern default, and that post-norm historically required careful warmup.
    \item \badgeLsix{L6} Derives the identity term in the pre-norm Jacobian, explains post-norm's depth-growing gradient imbalance at initialization (Xiong et al.), and states pre-norm's residual-stream-growth caveat.
    \item \badgeLseven{L7} Quantifies the stream growth ($\sqrt{l}$ at init), explains the ``effectively shallower'' consequence and its interaction with precision and attention logits, and maps each hybrid (sandwich, OLMo-2 reordering, peri-LN) to the specific failure it addresses; mentions DeepNet-style rescaling as the alternative route to deep post-norm.
\end{itemize}

\followups{``Why do modern recipes still use warmup if pre-norm fixes the gradients?'' ``What does the final norm before the LM head do in a pre-norm model?'' ``How does norm placement interact with QK-norm and attention-logit growth?''}
\end{interviewq}

\subsection{QK-Norm and Attention-Logit Growth}
\label{subsec:qknorm}

The $1/\sqrt{d_k}$ attention scaling (Chapter~\ref{chap:attention_transformers}) controls logit variance \textit{at initialization}: for random query/key vectors with unit-variance entries, the dot product has standard deviation $\sqrt{d_k}$, and dividing by $\sqrt{d_k}$ restores $O(1)$ logits. But the scaling is a constant. Training is free to grow the projections, and in large models it does: query/key norms increase and queries co-align with keys, so the logits $\vq^\top \vk / \sqrt{d_k}$ grow without bound over the run.

\textbf{The failure mode: entropy collapse.} Work the arithmetic for one head with $d_h = 128$. At initialization, logits are $O(1)$. Suppose training grows $\|\vq\|$ and $\|\vk\|$ to $10\sqrt{d_h} \approx 113$ each (a $10\times$ norm growth) and aligns them to cosine similarity $0.3$. The logit becomes $\|\vq\|\|\vk\|\cos\theta/\sqrt{d_h} = 113 \times 113 \times 0.3 / 11.3 \approx 340$. Once the gap between the top logit and the runner-up exceeds even $\sim$20, softmax is numerically one-hot (weight $1 - e^{-20}$), the attention distribution's entropy collapses to zero, and the softmax gradient (proportional to $p(1-p)$) vanishes---the head freezes. Precision makes it worse: BF16 has an 8-bit significand, so at logit magnitude 340 the quantization spacing is $\approx 2$---logit differences smaller than that no longer exist. This is exactly what killed ViT-22B's early training runs: within a few thousand steps, uncontrolled attention-logit growth (orders of magnitude beyond the $O(1)$ init scale) produced near one-hot attention weights and then divergence. The fix---normalizing queries and keys before the dot product---trained stably at 22B and became standard practice.

\begin{mathresult}{QK-Norm}
Apply RMSNorm to queries and keys, per head, before the dot product:
\begin{equation}
z_{ij} = \frac{\mathrm{RMSNorm}_{\gamma_q}(\vq_i) \cdot \mathrm{RMSNorm}_{\gamma_k}(\vk_j)}{\sqrt{d_h}}
\end{equation}
Each normalized vector has RMS 1, hence norm $\sqrt{d_h}$ (times its learned gain), so by Cauchy--Schwarz
\begin{equation}
|z_{ij}| \leq \sqrt{d_h}\,\gamma_q \gamma_k
\end{equation}
The attention logits are bounded \emph{regardless of what the projections learn}---for $d_h = 128$ and unit gains, $|z| \leq 11.3$.
\end{mathresult}

The cost is two extra RMSNorms over $d_h$-dimensional vectors per token---negligible against the attention matmuls. ViT-22B used the LayerNorm form; Gemma~3, Qwen3, and OLMo~2 apply RMSNorm to Q and K, and it is the expected first answer to ``how do you keep a 30B+ model from diverging.''

\textbf{The alternative: logit soft-capping (Gemma~2)---and its removal.} Gemma~2 instead bounded logits with a smooth squash, $z \leftarrow c \cdot \tanh(z/c)$, with $c = 50$ for attention logits and $c = 30$ for final output logits. This caps the magnitude, but as logits approach the cap the $\tanh$ compresses their differences and its gradient $1 - (z/c)^2$ shrinks---and the nonstandard softmax broke FlashAttention kernel compatibility when Gemma~2 launched. Gemma~3 dropped soft-capping in favor of QK-norm; capping now survives mainly as a mid-run retrofit when a spiking model cannot be re-architected.

\begin{keyinsight}
The $1/\sqrt{d_k}$ scaling and QK-norm solve two different problems. The scaling is a constant chosen for the \emph{initialization-time} statistics of random projections; it cannot respond to anything training learns. QK-norm constrains the \emph{learned} query/key magnitudes at every step, bounding the logits for the whole run---which is why models using QK-norm still keep the $1/\sqrt{d_h}$ factor. On large runs, the max attention logit per layer is the standard early-warning metric: it climbs for thousands of steps before the first loss spike.
\end{keyinsight}

\subsection{Normalization Selection Guide}

\begin{table}[H]
\centering
\caption{Normalization selection by architecture}
\begin{tabularx}{\textwidth}{lXX}
\toprule
\textbf{Architecture} & \textbf{Recommended} & \textbf{Why} \\
\midrule
CNN (large batch) & BatchNorm & Proven, effective \\
CNN (small batch) & GroupNorm & Batch-independent \\
Transformer & LayerNorm & Sequence-friendly \\
Modern LLM & RMSNorm & Simpler, similar quality; speedup is op-level only ($\sim$1\% end-to-end) \\
Style transfer & InstanceNorm & Per-image normalization \\
GAN discriminator & SpectralNorm & Stabilizes training \\
\bottomrule
\end{tabularx}
\end{table}

\begin{warningbox}
BatchNorm in multi-GPU training requires explicit synchronization (SyncBatchNorm). Without it, each GPU computes statistics from its local batch, which can be too small and lead to noisy estimates. This is a common source of subtle training bugs when scaling to multiple GPUs---loss looks fine but model quality degrades.
\end{warningbox}

%--- Q2: NaN training loss debugging (after normalization content) ---

\begin{interviewq}
\textbf{Q: Training loss is NaN after a few hundred steps. Which activation/normalization issues could cause this?}

\badgeLfive{L5} \quad \badgetype{Debugging} \quad \badgefreq{\freqhigh}

\stronganswer

NaN loss within the first few hundred steps points to numerical instability rather than slow divergence. I would diagnose in order of likelihood.

\textbf{Check 1: Gradient explosion before NaN.} Insert gradient norm logging. If gradient norms spike (e.g., from $\sim$1 to $>10^4$) before the NaN step, the issue is gradient explosion. Common causes: (a) post-norm placement---gradients must flow through normalization layers, which can amplify them in deep networks; switching to pre-norm often fixes this immediately; (b) no gradient clipping---add max-norm clipping at 1.0; (c) learning rate too high for the architecture depth.

\textbf{Check 2: Activation overflow in softmax or attention.} Large pre-softmax logits cause $\exp(\cdot)$ overflow. In attention, if Q/K projections are poorly initialized, dot products can be $O(d)$ instead of $O(1)$, making logits huge. Verify that attention is scaled by $1/\sqrt{d_k}$. Also check that temperature/scaling parameters haven't been accidentally set to values near zero (which amplifies logits).

\textbf{Check 3: Division by near-zero in normalization.} BatchNorm or LayerNorm compute $\hat{x} = (x - \mu)/\sqrt{\sigma^2 + \epsilon}$. If $\epsilon$ is too small (e.g., $10^{-12}$ instead of $10^{-5}$) and a batch has near-constant activations, the denominator approaches zero. This is especially dangerous in mixed-precision training where FP16 has limited range---$\epsilon$ should be at least $10^{-5}$ for FP16.

\textbf{Check 4: Dying ReLU cascade.} With aggressive initialization or high learning rates, a large fraction of ReLU neurons can die (output zero for all inputs). This creates zero gradients for those neurons, concentrating gradient flow through fewer active neurons, which then receive larger updates and can overflow. Check the fraction of dead neurons per layer; if $>$50\%, switch to GELU or LeakyReLU.

\textbf{Check 5: Mixed-precision loss scaling.} If using FP16 training, the loss scaler may not be configured correctly. FP16 underflows at $\sim 6 \times 10^{-8}$ and overflows at $\sim 65504$. Without dynamic loss scaling, gradients can underflow to zero (then weights don't update, then loss spikes) or overflow to Inf/NaN. Verify that the GradScaler is enabled and its initial scale is appropriate.

\redflags
\begin{itemize}
    \item ``Just lower the learning rate''---this is a symptom-level fix without diagnosing the root cause.
    \item Not mentioning mixed-precision training issues---this is the most common NaN source in modern training.
    \item Forgetting to check $\epsilon$ values in normalization layers.
\end{itemize}

\interviewtest{Systematic debugging approach for numerical instability, understanding the chain of events that leads to NaN (overflow $\to$ Inf $\to$ NaN in a subsequent operation), and awareness of the interaction between activations, normalization, and precision.}

\goodgreat
\begin{itemize}
    \item \badgeLfive{L5} Identifies the top 3 causes (softmax overflow, gradient explosion, normalization epsilon) and proposes concrete fixes.
    \item \badgeLsix{L6} Adds mixed-precision analysis, explains the FP16 range limitations quantitatively, and proposes systematic instrumentation (gradient norms, activation histograms, per-layer NaN detection) to localize the failing layer.
    \item \badgeLseven{L7} Discusses the interaction between initialization, normalization, and gradient flow in deep networks; mentions that pre-norm architectures are inherently more stable because the residual stream has an uninterrupted gradient path; references the ``edge of stability'' phenomenon where training can oscillate near the stability threshold before diverging.
\end{itemize}

\followups{``How do you instrument a training run to catch NaN before it propagates?'' ``What is the relationship between normalization placement and gradient magnitude?'' ``How does BF16 compare to FP16 for training stability?''}
\end{interviewq}

%--- Q8: Large-scale loss-spike debugging (after Q2; replaces the former validation-gap question, merged into Q5) ---

\begin{interviewq}
\textbf{Q: Your 70B-parameter run in BF16 starts loss-spiking at 100B tokens. Walk me through the activation- and normalization-level causes and mitigations.}

\badgeLsix{L6} \badgeLseven{L7} \quad \badgetype{Debugging} \quad \badgefreq{\freqhigh}

\stronganswer

First, frame the problem correctly: a run that was healthy for 100B tokens does not have a code bug---bugs fire immediately. Late-onset spikes are \textbf{slow-growth instabilities}: some quantity has been drifting upward for the whole run and has now crossed a threshold, usually interacting with BF16's coarse precision (8-bit significand, relative spacing $\approx 0.4\%$; range is not the issue---BF16 shares FP32's exponent). Also distinguish severity: spikes that recover on their own are warnings; a spike that does not recover is divergence. The same causes produce both.

\textbf{Cause 1: Attention-logit growth $\to$ entropy collapse.} Learned query/key norms and alignment grow over the run until softmax saturates (Section~\ref{subsec:qknorm}). The tell is in the logs \textit{before} the spike: max attention logit per layer/head climbs steadily for thousands of steps and then some layer crosses into one-hot attention. If this metric is not being logged, adding it is step one---it is the single best early-warning signal. Mitigation: QK-norm (RMSNorm on Q and K per head) for the next run or at a restart; logit soft-capping ($c\cdot\tanh(z/c)$) as a mid-run retrofit that does not change parameter shapes.

\textbf{Cause 2: Output-logit divergence at the final softmax.} The log-partition $\log Z$ of the output softmax drifts away from zero, softmax gradients degrade, and the cross-entropy becomes spike-prone. The standard mitigation is z-loss ($\sim 10^{-4}\log^2 Z$, PaLM), which regularizes $\log Z$ toward zero at negligible cost---mechanics in Chapter~\ref{chap:loss_functions}. Log $\log Z$ alongside max attention logits.

\textbf{Cause 3: Embedding-norm growth.} Embedding rows---especially for frequent tokens---grow throughout training, and with tied or large unembeddings this feeds directly into logit magnitude and into the scale of the layer-0 residual stream. Monitor embedding-norm statistics; mitigations include weight decay applied to embeddings and normalizing the embedding output (as OLMo does) so downstream layers see a controlled scale.

\textbf{Cause 4: Normalization epsilon and precision of the norm ops.} As residual-stream RMS grows over the run, norm layers divide by ever-larger statistics computed from ever-coarser BF16 values. Two checks: (a) norms (and softmax) should be computed in FP32 even in a BF16 run---a norm computed natively in BF16 loses the small differences it is supposed to standardize; (b) $\epsilon$ must suit the precision and the data---too small ($10^{-12}$) and a degenerate, near-constant token (e.g., repeated padding or BOS-heavy batches) produces a huge post-norm activation; RMSNorm's $\epsilon$ sits inside the square root (Section~\ref{sec:normalization}), and $10^{-6}$--$10^{-5}$ is the safe band.

\textbf{Cause 5: It may be the data, not the model.} Before touching architecture, test reproducibility: rewind to the checkpoint before the spike and replay the same batches. If the spike reproduces, inspect those batches (corrupt documents, pathological repetition); the pragmatic mitigation is the skip-batch strategy---rewind and skip a few hundred data batches (PaLM's recipe)---covered with the optimizer-side story (e.g., Adam $\beta_2 = 0.95$, gradient clipping) in Chapter~\ref{chap:training_optimization}. If the spike does \textit{not} reproduce, suspect hardware (a flaky node emitting bad gradients).

\textbf{The monitoring checklist} that separates candidates who have run this from those who have read about it: per-layer gradient norm, max attention logit, $\log Z$, residual-stream RMS per layer, embedding norms. All are nearly free to log, and every cause above announces itself in one of them thousands of steps before the first spike.

\redflags
\begin{itemize}
    \item ``Lower the learning rate''---a blunt instrument that trades away compute efficiency without identifying which quantity is growing; the spike usually returns.
    \item Reaching for FP16-era fixes (loss scaling, GradScaler)---BF16 does not underflow the way FP16 does; the issue is precision granularity, not range.
    \item Treating this as the single-GPU NaN checklist---at 100B tokens the causes are slow-growth (logits, embeddings, residual RMS), not initialization or a wrong $\epsilon$ that would have failed at step 100.
\end{itemize}

\interviewtest{Whether you have actually operated a large training run: the bug-vs-instability distinction, the monitoring-first mindset, knowledge of the specific quantities that grow (attention logits, $\log Z$, embedding and residual norms), and matching each mitigation to its mechanism rather than reciting a fix list.}

\goodgreat
\begin{itemize}
    \item \badgeLfive{L5} Proposes gradient-norm logging, checkpoint rewind, and checks data quality; knows BF16 vs.\ FP16 basics.
    \item \badgeLsix{L6} Distinguishes slow-growth instabilities from bugs, names attention-logit growth and QK-norm, knows z-loss and the skip-batch strategy, and insists norms/softmax run in FP32.
    \item \badgeLseven{L7} Lays out the full monitoring suite and the order in which signals fire; discusses trade-offs among mitigations (QK-norm vs.\ soft-capping retrofit, embedding weight decay vs.\ embedding norm); connects spike frequency to scale (larger models sit closer to the edge of stability) and knows the published lineage (PaLM's rewind-and-skip, ViT-22B's QK-norm, Gemma~2's capping and its removal in Gemma~3).
\end{itemize}

\followups{``Why do loss spikes become more frequent as models scale?'' ``The spike recovered on its own---do you intervene anyway?'' ``How would $\mu$P or a different parameterization change this failure surface?''}
\end{interviewq}
